\section{Main Theorem}

\begin{theorem}[Affine coordinate-pivot sweep and exact specializations]
\label{thm:main}
Fix $N\geq1$, $R>0$, $\kappa\in(0,\infty)$, a compact interval
$\Theta\subset\mathbb R$, and a deterministic affine family $(b,F)$.
Under Assumptions~\ref{assump:shared-pfaffian-chain} and
\ref{assump:no-forced-root},
\[
0\leq\Gamma_{\rm piv}(b,F;R)<\infty.
\]
Under Assumption~\ref{assump:joint-density-cap} as well, every
$\mu\in\mathcal D_{N,R,\kappa}$ and every interval $I\subseteq\Theta$
with $|I|>0$ satisfy
\[
\Pr_{\alpha\sim\mu}\!\left[\exists\theta\in I:
 b(\theta)+\langle\alpha,F(\theta)\rangle=0\right]
\leq
\kappa(2R)^{N-1}\Gamma_{\rm piv}(b,F;R)|I|
=\frac{A\Gamma_{\rm piv}(b,F;R)}{2R}|I|.
\tag{T1}\label{eq:t1}
\]
With the interval supremum taken before the law supremum and with
$\sup\varnothing=-\infty$,
\[
\sup_{\mu\in\mathcal D_{N,R,\kappa}}
\sup_{\substack{I\subseteq\Theta,\ I\text{ an interval}\\|I|>0}}
\frac{\Pr_{\alpha\sim\mu}[\exists\theta\in I:
 b(\theta)+\langle\alpha,F(\theta)\rangle=0]}{|I|}
\leq\frac{A\Gamma_{\rm piv}(b,F;R)}{2R}.
\tag{T2}\label{eq:t2}
\]
The left side of \eqref{eq:t2} is a finite nonnegative real precisely in
the branch in which both index classes are nonempty, namely
$A\geq1$ and $|\Theta|>0$.  If $A<1$ or $|\Theta|=0$, it is literally
$-\infty$; in that branch \eqref{eq:t2} is a vacuous extended-real
inequality and is not a finite nonnegative capacity statement.

For every $0<\delta\leq1$, including $\delta=1$, the scale-stress family
satisfies the exact identity
\[
\Gamma_{\rm piv}(b_\delta,F_\delta;1)=\frac1\delta.
\tag{T3}\label{eq:t3}
\]

For every integer $d\geq1$, every $R>0$, and every compact restriction of
the exact monic family,
\[
\Gamma_{\rm piv}(b_d,F_d;R)
\leq d+\frac{R d(d-1)}2.
\tag{T4}\label{eq:t4}
\]
Consequently, if $\mu\in\mathcal D_{d,R,\kappa}$ is any full joint law of
the $d$ lower coefficients, with arbitrary correlation, then every bounded
interval $I\subset\mathbb R$, including $|I|=0$, satisfies
\[
\Pr_{\alpha\sim\mu}\!\left[\exists\theta\in I:p_\alpha(\theta)=0\right]
\leq
\kappa(2R)^{d-1}
\left(d+\frac{R d(d-1)}2\right)|I|.
\tag{T5}\label{eq:t5}
\]
The leading coefficient one in \eqref{eq:t5} is deterministic and is not
part of the random vector.

All displayed rates have no hidden constants.  Their exposed variables are
$N,R,\kappa,A,\Gamma_{\rm piv},|I|$ and, in the specializations,
$\delta$ or $d$; the Pfaffian descriptors are fixed family descriptors and
do not enter as unshown factors.  The probabilities are ordinary
probabilities for each fixed law, the interval statements are static and
uniform over their stated interval classes, and there is no confidence or
horizon parameter.  The norm and measure conventions are the $\ell_1$
feasibility test defining $K_R$, scalar absolute coordinate-ratio variation,
one-dimensional interval length, and full-dimensional Lebesgue density.
No independence or probability conversion is implicit.  The theorem does
not bound $\Gamma_{\rm piv}$ polynomially in general Pfaffian presentation
data.
\end{theorem}
